\documentclass[11pt]{article}
\usepackage[hmargin=2.5cm,vmargin=2.5cm]{geometry}
\usepackage{amsmath,amssymb}
\usepackage[colorlinks=true,linkcolor=blue,urlcolor=blue]{hyperref}

\title{The Pricing Equations:\\ Calvo Price Setting with Indexation}
\author{Spending-efficiency model --- internal note}
\date{\today}

\begin{document}
\maketitle

\noindent
This note derives the four pricing equations of the model --- the two auxiliary
recursions for $x_{1,t}$ and $x_{2,t}$, the reset-price condition
$\epsilon x_{1,t} = (\epsilon-1)x_{2,t}$, and the price-index law of motion ---
from the retailer's problem (\texttt{models/model\_block.modpart}, price-setting
block). Notation is the paper's throughout.

\section{Setup}

Retailers purchase intermediate output at the relative price $mc_t$ --- their real
marginal cost --- differentiate it one-for-one into varieties, and sell variety $i$
at the nominal price $P_t(i)$. The final-good producer's cost minimization gives
the demand for each variety:
\begin{equation}\label{eq:demand}
    y_t(i) = \left(\frac{P_t(i)}{P_t}\right)^{-\epsilon} Y_t^{d}.
\end{equation}
Prices are sticky \`a la Calvo: each period a retailer reoptimizes its price with
probability $1-\theta_p$. With probability $\theta_p$ it cannot, and instead
indexes its price to past inflation with elasticity $\chi_p$, so a non-reoptimized
price is scaled by $\Pi_{t-1}^{\chi_p}$ each period.

Consider a retailer that reoptimizes at $t$. If it is still stuck with that price
at $t+s$, indexation will have scaled it by the cumulative factor
\begin{equation}\label{eq:index}
    X_{t,s} = \prod_{k=1}^{s} \Pi_{t+k-1}^{\chi_p}, \qquad X_{t,0} = 1.
\end{equation}

\section{The Reset-Price Problem}

The retailer chooses $P_t(i)$ to maximize the expected discounted real profits it
earns while the price remains in force. A profit $s$ periods ahead survives with
probability $\theta_p^{s}$ and is valued with the household's discount factor
$\beta^{s}\lambda_{t+s}/\lambda_t$; multiplying through by the constant
$\lambda_t$, the problem can be written:
\begin{equation*}
    \max_{P_t(i)} \; \mathbb{E}_t \sum_{s=0}^{\infty}
    \left(\beta\theta_p\right)^{s} \lambda_{t+s}
    \left[\frac{P_t(i)\,X_{t,s}}{P_{t+s}} - mc_{t+s}\right]
    \left(\frac{P_t(i)\,X_{t,s}}{P_{t+s}}\right)^{-\epsilon} Y_{t+s}^{d}.
\end{equation*}
This is why $\lambda$ appears inside the model's auxiliaries below: profits are
valued in utility units. The FOC with respect to $P_t(i)$ is:
\begin{equation*}
    \mathbb{E}_t \sum_{s=0}^{\infty} \left(\beta\theta_p\right)^{s} \lambda_{t+s}
    \left[(1-\epsilon)\,P_t(i)^{-\epsilon}
    \left(\frac{X_{t,s}}{P_{t+s}}\right)^{1-\epsilon}
    + \epsilon\, mc_{t+s}\,P_t(i)^{-\epsilon-1}
    \left(\frac{X_{t,s}}{P_{t+s}}\right)^{-\epsilon}\right] Y_{t+s}^{d} = 0.
\end{equation*}
Multiplying by $P_t(i)^{\epsilon+1}$ and rearranging:
\begin{equation}\label{eq:focP}
    (\epsilon-1)\,P_t(i)\;\mathbb{E}_t \sum_{s=0}^{\infty}
    \left(\beta\theta_p\right)^{s} \lambda_{t+s}
    \left(\frac{X_{t,s}}{P_{t+s}}\right)^{1-\epsilon} Y_{t+s}^{d}
    = \epsilon\;\mathbb{E}_t \sum_{s=0}^{\infty}
    \left(\beta\theta_p\right)^{s} \lambda_{t+s}\, mc_{t+s}
    \left(\frac{X_{t,s}}{P_{t+s}}\right)^{-\epsilon} Y_{t+s}^{d}.
\end{equation}
Nothing in \eqref{eq:focP} depends on $i$, so all reoptimizing retailers choose the
same reset price; call it $P_t^{\#}$.

\section{The Recursive Form}

\eqref{eq:focP} contains infinite sums and the price level; to obtain the model's
stationary system, express everything in relative terms. Define the relative reset
price and the indexation-deflated price drift:
\begin{equation}\label{eq:defs}
    \Pi_t^{*} \equiv \frac{P_t^{\#}}{P_t}, \qquad
    D_{t,s} \equiv \frac{X_{t,s}\,P_t}{P_{t+s}}, \qquad D_{t,0} = 1.
\end{equation}
$D_{t,s}$ is the relative price at $t+s$ of a retailer that reset to the aggregate
price level at $t$ and indexed ever since. Splitting off the first indexation step,
it obeys:
\begin{equation}\label{eq:Drec}
    D_{t,s} = \frac{\Pi_t^{\chi_p}}{\Pi_{t+1}}\; D_{t+1,s-1}.
\end{equation}
Now define the two sums in \eqref{eq:focP} in relative terms:
\begin{equation}\label{eq:x1def}
    x_{1,t} \equiv \mathbb{E}_t \sum_{s=0}^{\infty}
    \left(\beta\theta_p\right)^{s} \lambda_{t+s}\, mc_{t+s}\, D_{t,s}^{-\epsilon}\,
    Y_{t+s}^{d}, \qquad
    f_{t} \equiv \mathbb{E}_t \sum_{s=0}^{\infty}
    \left(\beta\theta_p\right)^{s} \lambda_{t+s}\, D_{t,s}^{1-\epsilon}\,
    Y_{t+s}^{d}.
\end{equation}
$x_{1,t}$ discounts marginal cost and $f_t$ discounts revenue, both weighted by the
demand the price will face. Using \eqref{eq:Drec}, splitting off the $s=0$ term
turns each sum into a recursion:
\begin{equation}\label{eq:x1}
    x_{1,t} = \lambda_t\, mc_t\, Y_t^{d}
    + \beta\theta_p\, \mathbb{E}_t\!\left[
    \left(\frac{\Pi_t^{\chi_p}}{\Pi_{t+1}}\right)^{-\epsilon} x_{1,t+1}\right],
\end{equation}
\begin{equation}\label{eq:f}
    f_{t} = \lambda_t\, Y_t^{d}
    + \beta\theta_p\, \mathbb{E}_t\!\left[
    \left(\frac{\Pi_t^{\chi_p}}{\Pi_{t+1}}\right)^{1-\epsilon} f_{t+1}\right].
\end{equation}
The model carries the revenue side as $x_{2,t} \equiv \Pi_t^{*} f_t$. Multiplying
\eqref{eq:f} by $\Pi_t^{*}$:
\begin{equation}\label{eq:x2}
    x_{2,t} = \lambda_t\, \Pi_t^{*}\, Y_t^{d}
    + \beta\theta_p\, \mathbb{E}_t\!\left[
    \left(\frac{\Pi_t^{\chi_p}}{\Pi_{t+1}}\right)^{1-\epsilon}
    \frac{\Pi_t^{*}}{\Pi_{t+1}^{*}}\, x_{2,t+1}\right].
\end{equation}
Finally, dividing \eqref{eq:focP} by $P_t$ and substituting the definitions, the
FOC becomes:
\begin{equation}\label{eq:reset}
    \epsilon\, x_{1,t} = (\epsilon-1)\, x_{2,t}.
\end{equation}
\eqref{eq:x1}, \eqref{eq:x2}, and \eqref{eq:reset} are the model equations
\texttt{x1 = lambda*mc*yd+...}, \texttt{x2 = lambda*PIstar*yd+...}, and
\texttt{epsilon*x1 = (epsilon-1)*x2}. \eqref{eq:reset} says the reset price is a
markup over (discounted) marginal cost: absent stickiness
($\theta_p=0$) it collapses to
$\Pi_t^{*} = \tfrac{\epsilon}{\epsilon-1}\,mc_t$, the static markup rule.

\section{The Price-Index Law of Motion}

The price index aggregates $P_t^{1-\epsilon} = \int_0^1 P_t(i)^{1-\epsilon}di$.
Each period a fraction $\theta_p$ of prices are last period's scaled by
$\Pi_{t-1}^{\chi_p}$ --- by the Calvo assumption the non-reoptimizers are a random
sample, so their price distribution is last period's --- and a fraction
$1-\theta_p$ equal $P_t^{\#}$:
\begin{equation*}
    P_t^{1-\epsilon}
    = \theta_p \left(\Pi_{t-1}^{\chi_p}\right)^{1-\epsilon} P_{t-1}^{1-\epsilon}
    + \left(1-\theta_p\right) \left(P_t^{\#}\right)^{1-\epsilon}.
\end{equation*}
Dividing by $P_t^{1-\epsilon}$:
\begin{equation}\label{eq:pindex}
    1 = \theta_p \left(\frac{\Pi_{t-1}^{\chi_p}}{\Pi_t}\right)^{1-\epsilon}
    + \left(1-\theta_p\right)\left(\Pi_t^{*}\right)^{1-\epsilon},
\end{equation}
the model equation \texttt{1 = thetap*(PI(-1)\^{}chi/PI)\^{}(1-epsilon)+
(1-thetap)*PIstar\^{}(1-epsilon)}. \eqref{eq:pindex} determines inflation: given
the reset price from \eqref{eq:reset}, it pins down $\Pi_t$.

\subsection{Aside: Price Dispersion Is an Aggregation Object}

The same reoptimizers/indexers split, applied to the demand weights
$\left(P_t(i)/P_t\right)^{-\epsilon}$ instead of $P_t(i)^{1-\epsilon}$, gives the
law of motion of price dispersion $v_t^{p} \equiv \int_0^1
\left(P_t(i)/P_t\right)^{-\epsilon} di$:
\begin{equation}\label{eq:vp}
    v_t^{p} = \theta_p \left(\frac{\Pi_{t-1}^{\chi_p}}{\Pi_t}\right)^{-\epsilon}
    v_{t-1}^{p} + \left(1-\theta_p\right)\left(\Pi_t^{*}\right)^{-\epsilon}.
\end{equation}
We derive it here because the algebra is one more application of the same split,
but $v_t^{p}$ is \emph{not} part of the pricing block: no retailer chooses
anything with respect to it. It is an aggregation statistic --- summing the variety
demands \eqref{eq:demand} gives $\int_0^1 y_t(i)\,di = v_t^{p}\,Y_t^{d} = Y_t$,
the wedge between gross output and demand --- and accordingly the paper and the
glossary keep it with market clearing.

\subsection{Aside: Why No Growth Factor in the Recursions}

Unlike the Euler equations, the recursions \eqref{eq:x1} and \eqref{eq:x2} carry no
growth factor in the model. The reason is that $\lambda$ and $Y^{d}$ enter only as
the product $\lambda_{t+s} Y_{t+s}^{d}$: the marginal utility scales inversely
with the trend and demand proportionally to it, so the product is stationary and
the detrending factors cancel term by term (see
\texttt{../detrending/detrending.tex} for the general logic).

\end{document}
